\documentclass[11pt,a4paper]{article}

\usepackage[margin=2.0cm]{geometry}
\usepackage[T1]{fontenc}
\usepackage{tgpagella}
\usepackage[scale=0.92]{tgheros}
\usepackage{amsmath,amssymb}
\usepackage{booktabs}
\usepackage{tabularx}
\usepackage{array}
\usepackage{enumitem}
\usepackage{xcolor}
\usepackage{titlesec}
\usepackage{fancyhdr}
\usepackage{tcolorbox}
\usepackage{hyperref}
\tcbuselibrary{skins,breakable}

\definecolor{accent}{HTML}{0E7490}
\definecolor{ink}{HTML}{1E293B}
\definecolor{soft}{HTML}{F0FDFA}
\definecolor{warnbg}{HTML}{FFF7ED}
\definecolor{ok}{HTML}{166534}
\definecolor{muted}{HTML}{64748B}
\definecolor{eqbg}{HTML}{F8FAFC}

\hypersetup{colorlinks=true,linkcolor=accent,urlcolor=accent}

\titleformat{\section}{\Large\sffamily\bfseries\color{accent}}{\thesection}{0.7em}{}
\titleformat{\subsection}{\large\sffamily\bfseries\color{ink}}{\thesubsection}{0.55em}{}
\titlespacing*{\section}{0pt}{12pt}{5pt}
\titlespacing*{\subsection}{0pt}{8pt}{3pt}

\setlist{itemsep=1.5pt,topsep=2pt,leftmargin=1.15em}
\setlength{\parskip}{2.5pt}
\renewcommand{\arraystretch}{1.18}
\newcolumntype{L}{>{\raggedright\arraybackslash}X}

\pagestyle{fancy}
\fancyhf{}
\fancyhead[L]{\small\sffamily\color{muted} ReactorGenome · Explain everything}
\fancyhead[R]{\small\sffamily\color{muted} CHE331 viva}
\fancyfoot[C]{\small\thepage}
\renewcommand{\headrulewidth}{0.3pt}

\tcbset{boxrule=0pt,fonttitle=\sffamily\bfseries,coltitle=white,breakable,sharp corners=southwest}
\newtcolorbox{pitchbox}{colback=soft,leftrule=3pt,colframe=accent,title={Pitch},colbacktitle=accent}
\newtcolorbox{eqbox}{colback=eqbg,leftrule=3pt,colframe=accent!70,boxrule=0pt,before skip=4pt,after skip=4pt}
\newtcolorbox{tipbox}{colback=warnbg,leftrule=3pt,colframe=orange!70,title={Say this},colbacktitle=orange!75!black}
\newtcolorbox{qabox}[1]{colback=white,colframe=accent!35,boxrule=0.35pt,title={Q: #1},colbacktitle=accent!90!black,before skip=6pt,after skip=4pt}

\newcommand{\term}[1]{\textbf{#1}}

\begin{document}

\begin{center}
{\sffamily\bfseries\fontsize{24}{28}\selectfont\color{accent} Explain the project to your professor\par}
\vspace{0.2cm}
{\large\sffamily How it was made · Maths · Graphs · Terms · Books\par}
\vspace{0.35cm}
{\small App: \url{https://che331.streamlit.app}
\quad·\quad
Report: \texttt{docs/report/report.pdf}
\quad·\quad
Tests: \href{https://github.com/Mayurdp2403/ReactorGenome/tree/main/tests}{GitHub / tests}}
\end{center}

\begin{pitchbox}
CHE331's four reactor labs are run under fixed conditions and analysed later. ReactorGenome is a laptop web app beside the
rig: type the same readings, get live fits and diagnostics, switch assumptions, then use your fitted \(k\) in a Design
Studio the lab cannot build. No new sensors.
\end{pitchbox}

%=====================================================================
\section{How the app was made (basics --- no code deep dive)}

\subsection{What you actually built}
A \textbf{web application}: open a browser, no separate desktop install. One site serves Batch, CSTR, PFR, RTD and a
Design Studio. Live URL: \url{https://che331.streamlit.app}.

\subsection{Tools (if they ask ``what did you use?'')}
\begin{tabularx}{\textwidth}{@{}l L@{}}
\toprule
\textbf{Piece} & \textbf{Role in plain words} \\
\midrule
Python & Language the whole project is written in. \\
Streamlit & Turns Python into a clickable web page (inputs, tables, tabs). \\
NumPy / SciPy & Numbers, fitting curves, solving ODEs (reactor profiles, Arrhenius, RTD). \\
Pandas & Readings tables (rows of \(T\), flow, conductivity). \\
Plotly & Interactive graphs (zoom, hover, save PNG). \\
HTML/JS board & Drag-and-drop plant drawing in Design Studio. \\
Streamlit Cloud & Free hosting so anyone with the link can open the app. \\
\bottomrule
\end{tabularx}

\subsection{Architecture (one picture in words)}
\begin{itemize}
  \item \textbf{Engine:} the chemistry and reactor equations only (no buttons). CSTR, PFR, batch, RTD, flowsheet solver.
  \item \textbf{Pages:} what the student sees --- one page per experiment + Design Studio.
  \item Same layout every experiment: \textbf{left} assumptions, \textbf{middle} live reactor, \textbf{right} results,
  \textbf{below} readings table, \textbf{tabs} analysis and export.
\end{itemize}

\subsection{How a number reaches the screen}
Student types conductivity \(\to\) app converts to conversion \(X\) \(\to\) builds the design-plot coordinate \(y\) \(\to\)
fits a line to get \(k\) at that \(T\) \(\to\) Arrhenius across temperatures \(\to\) \(E_a\) and predictions. Every switch
on the left changes the model and redraws.

\subsection{How you know it is correct}
Analytic limits (ideal CSTR/PFR, RTD moments, recycle) plus replaying 2026 lab-report sheets with their own conventions
(\(k\) and \(E_a\) match to four figures). Test suite: GitHub \texttt{tests} folder.

%=====================================================================
\section{Reference books and course material}

Say these names if asked ``where do the models come from?''

\begin{tabularx}{\textwidth}{@{}l L@{}}
\toprule
\textbf{Source} & \textbf{What you took from it} \\
\midrule
H.\,S.\ Fogler, \emph{Elements of Chemical Reaction Engineering} &
Mole balances, CSTR/PFR design equations, Arrhenius, non-ideal reactors, energy balance / multiple steady states,
Levenspiel plots, RTD (E, F, \(t_m\), \(\sigma^2\), tanks-in-series, dispersion / Pe). \\
O.\ Levenspiel, \emph{Chemical Reaction Engineering} &
Same reactor types; classic \(X\)--\(Da\) comparison; Levenspiel plot (area = size). \\
CHE331 lab manuals (Batch, CSTR, PFR, RTD) &
Observation-table layout, design-plot conventions, typed \(\Lambda_\infty\), how groups report \(k\) and \(E_a\). \\
Ionic conductivity data (limiting molar conductivities) &
Default \(\Lambda_\infty \approx 0.366\,\Lambda_0\) (OH\(^-\) replaced by acetate). \\
\bottomrule
\end{tabularx}

\begin{tipbox}
``I did not invent new reactor theory. I implemented the Fogler / Levenspiel balances and the lab-manual fitting
workflow in software so students can use them \emph{during} the three-hour slot.''
\end{tipbox}

%=====================================================================
\section{Chemistry and the link to conductivity}

\textbf{Reaction (saponification):}
\[
\mathrm{NaOH + CH_3COOC_2H_5 \rightarrow CH_3COONa + C_2H_5OH}
\]
Usually treated as \(-r_A = k\,C_A C_B\). Equimolar feed \(\Rightarrow -r_A = k\,C_A^2\).

\textbf{Why measure conductivity?} NaOH (OH\(^-\)) conducts strongly. As it is consumed, conductivity \(\Lambda\) falls.
\[
X = \frac{\Lambda_0 - \Lambda}{\Lambda_0 - \Lambda_\infty}
\]
\(\Lambda_0\): feed / start. \(\Lambda_\infty\): complete conversion. Too-low \(\Lambda_\infty\) \textbf{inflates \(X\)} and
\textbf{shrinks \(k\)}. App default \(\approx 0.366\,\Lambda_0\); manuals often type 0.1 or 3.75\,mS/cm.

%=====================================================================
\section{Mathematical models (what to write on the board)}

\subsection{Symbols you must know}
\begin{tabularx}{\textwidth}{@{}ll L@{}}
\toprule
Symbol & Name & Meaning \\
\midrule
\(C_A, C_B\) & Concentrations & NaOH, ester (mol/L) \\
\(C_{A0}\) & Feed concentration & After mixing, at reactor inlet \\
\(X\) & Conversion & Fraction of NaOH reacted \\
\(\tau = V/Q\) & Residence / space time & Tank or coil volume \(\div\) volumetric flow \\
\(k\) & Rate constant & Units depend on order (2nd: L/mol·min) \\
\(E_a, A\) & Activation energy, pre-factor & Arrhenius: \(k = A\,e^{-E_a/RT}\) \\
\(\mathrm{Da}\) & Damköhler number & \(k\,C_{A0}^{n-1}\tau\) --- reaction speed vs flow speed \\
\(M = C_{B0}/C_{A0}\) & Stoichiometric ratio & 1 = equimolar \\
\(N\) & Tanks in series & From RTD: \(N = t_m^2/\sigma^2\) \\
\(\mathrm{Pe}\) & Péclet number & Convection \(\div\) axial dispersion \\
\bottomrule
\end{tabularx}

\subsection{Ideal CSTR (steady, isothermal, 2nd order, equimolar)}
Mole balance: \(F_{A0}X = (-r_A)V\) with \(-r_A = k C_{A0}^2 (1-X)^2\).
\begin{eqbox}
\[
\tau = \frac{X}{k\,C_{A0}(1-X)^2}
\quad\Rightarrow\quad
y = \frac{X}{(1-X)^2} = k\,C_{A0}\,\tau
\]
\end{eqbox}
Plot \(y\) vs \(\tau\): slope \(= k\,C_{A0}\) \(\Rightarrow\) \(k\). (If \(M\neq 1\), \(y\) uses \((1-X)(M-X)\) form.)

\subsection{Ideal PFR (plug flow, same kinetics)}
\begin{eqbox}
\[
\tau = \int_0^X \frac{\mathrm{d}X'}{-r_A(X')}
\quad\Rightarrow\quad
y = \frac{X}{1-X} = k\,C_{A0}\,\tau
\quad\text{(equimolar 2nd order)}
\]
\end{eqbox}
Along the coil the app integrates the ODE \(\mathrm{d}C_A/\mathrm{d}V = r_A/Q\).

\subsection{Batch}
Time from concentration: \(\mathrm{d}C_A/\mathrm{d}t = r_A\). Integral method: for 2nd order equimolar,
\(1/C_A - 1/C_{A0} = k t\). Differential method: plot \(\ln(-\mathrm{d}C_A/\mathrm{d}t)\) vs \(\ln C_A\); slope \(\approx\) order \(n\).

\subsection{Arrhenius}
\begin{eqbox}
\[
\ln k = \ln A - \frac{E_a}{R}\,\frac{1}{T}
\]
\end{eqbox}
Slope of \(\ln k\) vs \(1/T\) is \(-E_a/R\). App reports 95\% intervals.

\subsection{Non-ideal CSTR (speak the idea)}
\begin{itemize}
  \item \textbf{Bypass \(\beta\):} fraction of feed skips the tank \(\Rightarrow\) lower overall \(X\).
  \item \textbf{Dead volume \(\alpha\):} fraction of \(V\) inactive \(\Rightarrow\) smaller effective \(\tau\).
  \item \textbf{\(N\) tanks in series:} same total \(V\), approaches PFR as \(N\) grows \(\Rightarrow\) higher \(X\).
\end{itemize}

\subsection{Non-ideal PFR}
\begin{itemize}
  \item \textbf{Axial dispersion:} Fogler open/closed vessel; characterised by Pe. High Pe \(\to\) plug flow.
  \item \textbf{Laminar:} velocity profile across radius; conversion below ideal plug for the same \(\tau\).
  \item \textbf{Sample delay:} reaction continues in the beaker; looks like a large free intercept on \(\tau\) vs \(y\).
\end{itemize}

\subsection{RTD}
\begin{eqbox}
\[
E(t)=\frac{C(t)}{\int_0^\infty C\,\mathrm{d}t},\quad
t_m=\int t E\,\mathrm{d}t,\quad
\sigma^2=\int (t-t_m)^2 E\,\mathrm{d}t,\quad
N=\frac{t_m^2}{\sigma^2}
\]
\end{eqbox}
\(F(t)=\int_0^t E\,\mathrm{d}t\). Dimensionless \(E(\theta)\) with \(\theta=t/t_m\). Tail cut-off underestimates \(\sigma^2\)
and inflates \(N\).

\subsection{Energy balance (CSTR jacket)}
At steady state, heat generation from reaction equals heat removal by the jacket \(UA(T-T_a)\) (plus sensible heat of
flow). Plot generation and removal vs \(T\): one or three intersections \(\Rightarrow\) multiple steady states / runaway.

\subsection{Flowsheet (Design Studio)}
Each block uses the same rate law. Mixers and splitters conserve mass. Recycle: iterate until compositions stop changing.
Overall \(X\) from all feeds to product. Mass balance on Na and acetyl groups must close (\(<10^{-8}\) relative).

%=====================================================================
\section{What every graph does}

\subsection{CSTR / PFR Results tab}
\begin{tabularx}{\textwidth}{@{}l L@{}}
\toprule
\textbf{Graph} & \textbf{What it shows / why it is there} \\
\midrule
Design plot (one \(T\)) & \(y\) vs \(\tau\) (or flipped). Line through points \(\to k\) at that temperature. Core of the manual. \\
Design plot (all \(T\)) & Same axes, several temperatures. Steeper slope = larger \(k\). \\
Arrhenius & \(\ln k\) vs \(1/T\). Slope \(\to E_a\). Proves temperature dependence is consistent. \\
\(X\) vs total flow & Higher flow \(\Rightarrow\) smaller \(\tau\) \(\Rightarrow\) usually lower \(X\). \\
\(X\) vs \(T\) & Higher \(T\) \(\Rightarrow\) larger \(k\) \(\Rightarrow\) usually higher \(X\). \\
\(k\) from each run alone & Each point is \(k\) from one row. Flat vs \(\tau\) = consistent data; sloping = \(\Lambda_\infty\), delay, or unsteady. \\
Parity (model vs measured \(X\)) & Points on the diagonal = fit is good. \\
Residuals & Measured \(-\) model for each run; shows which runs disagree. \\
\bottomrule
\end{tabularx}

\subsection{Predictions panel}
\begin{itemize}
  \item \(X\) vs \(\tau\) at chosen \(T\): size the reactor mentally.
  \item \(k\) vs \(T\): Arrhenius curve with your fitted points.
  \item \(X\) vs \(C_{A0}\) for orders 1--3: all orders fit at the lab \(C_{A0}\); they split if concentration changes --- that
  is how you tell order apart.
\end{itemize}

\subsection{Assumption effects}
Bar chart of \(X\) under ideal vs bypass / dead volume / \(N\) tanks / stagnant / reversible. Answers ``what was that
assumption costing me?''

\subsection{Energy balance}
Heat generation vs removal vs \(T\); phase portrait of trajectories. Shows multiple steady states and runaway basin.

\subsection{PFR: conversion along coil; flow models}
\(X(z/L)\): builds along the tube. Flow-models plot: same \(k\), different hydrodynamics, points overlaid.

\subsection{Batch}
\(C_A(t)\); integral plot (straight if order right); differential log--log (slope \(\approx n\)); residuals; AIC table ranking
rate laws.

\subsection{RTD}
\(E(t)\) with \(t_m\) marked; overlay of all runs; \(E(\theta)\), \(F(t)\); model fits (tanks-in-series, dispersion) + residuals.

\subsection{Design Studio}
Profiles (\(C_A\), \(T\) vs cumulative \(V\)); Levenspiel \(1/(-r_A)\) vs \(X\) (area \(\propto\) size); sweep of one parameter;
\(k(T)\) from every rig on one plot + ideal \(X\)--\(Da\) curves.

%=====================================================================
\section{Glossary --- terms the professor may ask}

\begin{description}[style=nextline,leftmargin=0pt,font=\sffamily\bfseries\color{accent}]
  \item[Residence time \(\tau\)] \(V/Q\). Average time fluid would spend if the vessel were ideally swept.
  \item[Space time] Same as \(\tau\) for constant density liquid (Fogler).
  \item[Conversion \(X\)] Fraction of limiting reactant (NaOH) that has reacted.
  \item[Damköhler Da] Dimensionless rate \(\times\) time. Large Da \(\Rightarrow\) high \(X\) for a given reactor type.
  \item[Design plot] Linearised kinetics plot (\(y\) vs \(\tau\)) whose slope gives \(k\).
  \item[Arrhenius] How \(k\) depends on \(T\); gives \(E_a\).
  \item[Activation energy \(E_a\)] Barrier; higher \(E_a\) \(\Rightarrow\) \(k\) rises faster with \(T\).
  \item[Pre-exponential \(A\)] Arrhenius intercept factor (same units as \(k\)).
  \item[Equimolar] \(C_{A0}=C_{B0}\) (\(M=1\)).
  \item[Steady state] Properties no longer change with time. CSTR outlet conductivity flat.
  \item[Transient / start-up] Approach to steady state after start or flow change.
  \item[Bypass] Feed fraction that never enters the reacting volume.
  \item[Dead volume] Stagnant fraction of \(V\) that does not participate.
  \item[Tanks in series \(N\)] Ideal CSTRs chained; large \(N\) \(\approx\) PFR.
  \item[RTD] Residence time distribution --- how long tracer molecules stay.
  \item[\(E(t)\)] Exit age density; probability density of leaving at time \(t\).
  \item[\(F(t)\)] Cumulative fraction that has left by time \(t\).
  \item[Péclet Pe] Convective transport / axial dispersion. High Pe \(\approx\) plug flow.
  \item[Levenspiel plot] \(1/(-r_A)\) vs \(X\); area related to reactor volume needed.
  \item[Parity plot] Predicted vs measured; checks the fit visually.
  \item[Residual] Measured minus model; should be small and random.
  \item[AIC] Information criterion ranking models; lower is better (batch rate laws).
  \item[Sample delay] Extra reaction time after the PFR outlet (beaker).
  \item[Multiple steady states] More than one stable \(T\) (and \(X\)) for the same feeds/jacket.
  \item[Runaway] Trajectory that jumps to the hot steady state.
  \item[UA] Jacket heat-transfer coefficient \(\times\) area (cooling strength).
\end{description}

%=====================================================================
\section{Page-by-page: what each experiment does}

\subsection{CSTR}
Simulate / animate tank; type lab-sheet readings; extrapolate to true steady \(\Lambda\); fit \(k(T)\) and \(E_a\);
assumption switches; energy balance; export.

\subsection{PFR}
Coil profiles; average readings per flow; design plots; plug / dispersion / laminar / recycle / delay; Re and Dean from
coil geometry.

\subsection{Batch}
\(\kappa(t)\to C_A\); choose rate law; clock offset; AIC comparison; Arrhenius if several \(T\).

\subsection{RTD}
Multi-run tracer sheet; moments; \(N\) and Pe (several conventions); tail warning; model ranking; conversion implied by
the RTD.

\subsection{Design Studio}
Build series / recycle / bypass / two feeds / jacketed tank; share \(k\) from lab pages; streams, profiles, Levenspiel,
sweep, compare rigs; save/load JSON.

%=====================================================================
\section{3-hour lab use (students + TA)}

\begin{enumerate}
  \item Open \url{https://che331.streamlit.app}, name the project.
  \item Set \(V\), \(C_{A0}\), \(\Lambda_\infty\) mode to match the manual.
  \item Run the experiment; type readings when due (CSTR stopwatch helps).
  \item Check diagnostics: settled or not? \(k\) vs \(\tau\) flat?
  \item Change flow/\(T\); watch Arrhenius fill in.
  \item Toggle one assumption; see \(\Delta X\).
  \item Publish kinetics \(\to\) Design Studio \(\to\) try series/recycle.
  \item Export ZIP; Project \(\to\) Save JSON to share the session.
\end{enumerate}

%=====================================================================
\section{Numbers to remember}

\begin{center}\small
\begin{tabular}{@{}ll@{}}
\toprule
CSTR report match & \(E_a=25.07\)\,kJ/mol; \(k=3.16/5.85/7.44\) at 31/46/60\,°C \\
PFR report match & \(E_a=12.70\)\,kJ/mol; \(k=11.74/13.75/15.73\) at 40/50/60\,°C \\
Design Studio demo & CSTR\(\to\)PFR overall \(X\approx 0.545\) \\
Default \(\Lambda_\infty/\Lambda_0\) & \(\approx 0.366\) (ions) \\
\bottomrule
\end{tabular}
\end{center}

%=====================================================================
\section{Short Q\&A}

\begin{qabox}{Did you derive new kinetics?}
No --- Fogler/Levenspiel mole balances and the lab-manual design plots, implemented and checked against reports.
\end{qabox}

\begin{qabox}{Where does \(k\) come from?}
From the slope of the design plot at each \(T\), then Arrhenius for \(E_a\) and \(A\).
\end{qabox}

\begin{qabox}{What is Da?}
\(\mathrm{Da}=k C_{A0}^{n-1}\tau\). Large Da means reaction is fast compared with the time fluid spends in the reactor.
\end{qabox}

\begin{qabox}{Why can CSTR and PFR \(E_a\) differ between groups?}
Different \(\Lambda_\infty\), scatter, sample delay, unsteady averages. App matches each report's method; fair comparison
needs the same conventions (Compare rigs).
\end{qabox}

\begin{qabox}{Best insight from the tool?}
Low typed \(\Lambda_\infty\) and averaging before steady state explain weird numbers better than ``wrong kinetics.''
\end{qabox}

\vspace{0.5cm}
\begin{center}
\color{accent}\sffamily\bfseries
Know the pitch, the design-plot equation, Arrhenius, and what each Results graph is for.\\
That is enough to explain the whole project.
\end{center}

\end{document}
